% 例 3：蚂蚁沿圆柱侧面爬行最短路径
% 圆柱 r=2, h=6π，A 在底圆周；B 在上底圆周"正对面"（绕半圈）
\documentclass[tikz,border=4pt]{standalone}
\usepackage{ctex}
\usepackage{amsmath}
\usetikzlibrary{calc, arrows.meta}
\begin{document}
\begin{tikzpicture}[thick, scale=0.7]
  \pgfmathsetmacro{\r}{1.0}
  \pgfmathsetmacro{\h}{3.2}
  \pgfmathsetmacro{\yr}{0.3}
  % 圆柱
  \draw (-\r,0) -- (-\r,\h);
  \draw (\r,0) -- (\r,\h);
  \draw (0,\h) ellipse ({\r} and {\yr});
  \draw (-\r,0) arc (180:360:{\r} and {\yr});
  \draw[dashed,gray] (-\r,0) arc (180:0:{\r} and {\yr});
  % A 在底前方点，B 在上后方
  \coordinate (A) at (\r,0);
  \filldraw (A) circle (1.5pt);
  \node[below right, font=\footnotesize] at (A) {$A$};
  \coordinate (B) at (-\r,\h);
  \filldraw (B) circle (1.5pt);
  \node[above left, font=\footnotesize] at (B) {$B$};
  % 路径示意（弯曲）
  \draw[red, thick, dashed] (A) .. controls (\r, \h/2) and (-\r, \h/2-0.5) .. (B);
  \node[below, font=\footnotesize] at (0,-0.8) {圆柱表面};

  % 展开矩形：长 2πr=4π ≈ 6.3，宽 h=6π ≈ 6.3。画成长 4 宽 4.5 的矩形示意
  \begin{scope}[xshift=5cm, yshift=0cm]
    \pgfmathsetmacro{\L}{4}
    \pgfmathsetmacro{\H}{4.5}
    \draw (0,0) rectangle (\L, \H);
    \coordinate (Ap) at (0,0);
    \coordinate (Bp) at (\L/2, \H);  % A 处剪开后，对面在水平 πr = L/2
    \filldraw (Ap) circle (1.5pt);
    \filldraw (Bp) circle (1.5pt);
    \node[below left, font=\footnotesize] at (Ap) {$A$};
    \node[above, font=\footnotesize] at (Bp) {$B$};
    \draw[red, thick] (Ap) -- (Bp);
    % 直角边
    \draw[gray, dashed] (Ap) -- (\L/2, 0) -- (Bp);
    \node[font=\footnotesize, below, gray] at (\L/4, 0) {$\pi r=2\pi$};
    \node[font=\footnotesize, right, gray] at (\L/2, \H/2) {$h=6\pi$};
    \node[font=\footnotesize, left, red] at ($(Ap)!0.5!(Bp)$) {$AB$};
    \node[below, font=\footnotesize] at (\L/2, -0.8) {展开后用勾股};
  \end{scope}
\end{tikzpicture}
\end{document}
